\documentclass[12pt,a4paper]{ctexart}
\usepackage{amsmath,amssymb}
\usepackage{geometry}
\usepackage{listings}
\usepackage{xcolor}
\usepackage{hyperref}
\hypersetup{
	colorlinks=true,
	linkcolor={blue!60!black},
	urlcolor={blue!60!black},
	citecolor={blue!60!black},
	pdfborder={0 0 0},
}
\usepackage{tabularx}
\usepackage{fancyhdr}
\usepackage{tikz}
\usepackage{lastpage}
\usepackage{enumitem}
\usepackage{array}
\usepackage{booktabs}
\usepackage{longtable}
\geometry{left=2.5cm,right=2.5cm,top=2.5cm,bottom=2.5cm}
\definecolor{codebg}{RGB}{245,245,245}
\definecolor{codeframe}{RGB}{200,200,200}
\definecolor{commentcolor}{RGB}{0,128,0}
\definecolor{keywordcolor}{RGB}{127,0,85}
\definecolor{stringcolor}{RGB}{42,0,255}
\definecolor{accent}{HTML}{4F46E5}
\definecolor{accent2}{HTML}{7C3AED}
\definecolor{mid}{HTML}{6B7280}
\lstset{
	language=C++,
	basicstyle=\small\ttfamily,
	backgroundcolor=\color{codebg},
	frame=single,
	rulecolor=\color{codeframe},
	breaklines=true,
	showstringspaces=false,
	commentstyle=\color{commentcolor},
	keywordstyle=\color{keywordcolor}\bfseries,
	stringstyle=\color{stringcolor},
	numbers=left,
	numberstyle=\tiny\color{gray},
	tabsize=4,
	captionpos=b,
	extendedchars=true,
	columns=flexible,
}
\pagestyle{fancy}
\fancyhf{}
\fancyhead[L]{\leftmark}
\fancyhead[R]{并查集算法专题}
\fancyfoot[C]{\thepage\ /\ \pageref*{LastPage}}
\renewcommand{\headrulewidth}{0.4pt}
\renewcommand{\footrulewidth}{0pt}
\newcommand{\infobox}[1]{%
	\par\noindent
	\fcolorbox{accent!40}{accent!5}{%
		\begin{minipage}{\dimexpr\textwidth-2\fboxsep-2\fboxrule\relax}
			#1
		\end{minipage}%
	}\par
}
\title{并查集算法专题}
\author{算法学习笔记}
\date{\today}
\begin{document}
	\maketitle
	\tableofcontents
	\newpage
	\section{并查集基础}
	\subsection{标准并查集（路径压缩+按秩合并）}
	\infobox{解决问题：维护一个元素集合，支持两种操作：合并两个元素所在的集合、查询两个元素是否属于同一集合。适用于社交网络好友关系判断、连通分量统计、Kruskal最小生成树等场景。}
	\subsubsection*{题目描述}
	给定n个元素（编号1~n），初始每个元素自成一个集合。有m次操作，操作类型有两种：1 x y表示合并x和y所在的集合；2 x y表示查询x和y是否在同一个集合中。
	\subsubsection*{输入示例}
	\begin{verbatim}
		5 6
		1 1 2
		1 2 3
		2 1 3
		1 4 5
		2 1 4
		2 4 5
	\end{verbatim}
	\subsubsection*{输出示例}
	Yes
	No
	Yes
	\begin{lstlisting}
		#include <bits/stdc++.h>
		using namespace std;
		struct DSU {
			vector<int> parent;  // parent[i]: 节点i的父节点，根节点的父节点指向自己
			vector<int> rank;    // rank[i]: 节点i所在的树的高度（按秩合并时使用）
			DSU(int n) {  // 构造函数：初始化并查集
				parent.resize(n + 1);  // 索引从1到n，所以分配n+1个空间
				rank.resize(n + 1, 0); // 初始所有节点高度为0
				for (int i = 1; i <= n; i++) {
					parent[i] = i;  // 初始时每个节点的父节点都是自己，即每个元素自成一个集合
				}
			}
			// 核心原理：查找节点x所在集合的根节点（代表元），同时进行路径压缩优化
			// 路径压缩的作用：将查找路径上的所有节点直接指向根节点，降低后续查找的时间复杂度
			int find(int x) {
				if (parent[x] != x) {  // 如果x不是根节点
					parent[x] = find(parent[x]);  // 递归查找父节点，并将结果赋值给parent[x]（路径压缩）
				}
				return parent[x];  // 返回根节点
			}
			// 合并两个节点所在的集合
			void unite(int x, int y) {
				int rootX = find(x);  // 找到x的根节点
				int rootY = find(y);  // 找到y的根节点
				if (rootX == rootY) return;  // 如果已经在同一集合，无需合并
				// 按秩合并：将高度较小的树合并到高度较大的树，保持树平衡
				if (rank[rootX] < rank[rootY]) {
					parent[rootX] = rootY;  // 将rootX的父节点设为rootY
				} else if (rank[rootX] > rank[rootY]) {
					parent[rootY] = rootX;  // 将rootY的父节点设为rootX
				} else {
					parent[rootY] = rootX;  // 高度相等时，任意合并，并将高度+1
					rank[rootX]++;  // 新树的高度增加1
				}
			}
			// 查询两个节点是否在同一个集合中
			bool same(int x, int y) {
				return find(x) == find(y);
			}
		};
		int main() {
			int n, m;  // n: 元素总数, m: 操作次数
			cin >> n >> m;
			DSU dsu(n);  // 创建大小为n的并查集
			for (int i = 0; i < m; i++) {
				int op, x, y;  // op: 操作类型, x, y: 操作的两个元素
				cin >> op >> x >> y;
				if (op == 1) {
					dsu.unite(x, y);  // 合并操作
				} else {
					if (dsu.same(x, y)) {  // 查询操作
						cout << "Yes\n";
					} else {
						cout << "No\n";
					}
				}
			}
			return 0;
		}
	\end{lstlisting}
	\subsection{带权并查集（维护到根节点的距离）}
	\infobox{解决问题：维护元素之间的相对关系，如食物链中的捕食关系、向量偏移、区间和关系等。适用于解决带关系的集合合并问题，每个元素到根节点的距离有意义。}
	\subsubsection*{题目描述}
	有n个动物（编号1~n），被分成三类：A吃B，B吃C，C吃A。给出k句话，每句话描述x和y的关系（1表示x和y同类，2表示x吃y）。判断假话的数量（假话条件：与前面的话冲突，或者编号超出范围）。
	\subsubsection*{输入示例}
	\begin{verbatim}
		100 7
		1 101 1
		2 1 2
		2 2 3
		2 3 3
		1 1 3
		2 3 1
		1 5 5
	\end{verbatim}
	\subsubsection*{输出示例}
	3
	\begin{lstlisting}
		#include <bits/stdc++.h>
		using namespace std;
		struct WeightedDSU {
			vector<int> parent;  // parent[i]: 节点i的父节点
			vector<int> dist;    // dist[i]: 节点i到其父节点的关系（0表示同类，1表示被父节点吃，2表示吃父节点）
			WeightedDSU(int n) {  // 构造函数：初始化带权并查集
				parent.resize(n + 1);
				dist.resize(n + 1, 0);  // 初始到父节点的距离为0（自己到自己是同类）
				for (int i = 1; i <= n; i++) {
					parent[i] = i;  // 初始每个节点自成一棵树
				}
			}
			// 核心原理：查找根节点时，同时更新dist[x]为x到根节点的关系
			// 关系通过路径上的关系值相加并取模3得到（因为只有3种关系）
			int find(int x) {
				if (parent[x] != x) {
					int root = find(parent[x]);  // 递归找到根节点
					dist[x] = (dist[x] + dist[parent[x]]) % 3;  // 更新x到根节点的关系
					parent[x] = root;  // 路径压缩，直接指向根节点
				}
				return parent[x];
			}
			// 合并操作：根据给定关系type合并x和y
			// type=0表示x和y同类，type=1表示x吃y
			bool unite(int x, int y, int type) {
				int rootX = find(x);  // 找到x的根节点
				int rootY = find(y);  // 找到y的根节点
				if (rootX == rootY) {  // 已经同根，检查是否与已有关系矛盾
					int relation = (dist[x] - dist[y] + 3) % 3;  // 计算x到y的关系
					return relation == type;  // 如果计算出的关系与给定的type一致，则真话；否则假话
				}
				// 合并两个树：将rootX合并到rootY，并更新dist[rootX]
				parent[rootX] = rootY;
				// 推导dist[rootX]的公式：已知dist[x]到rootX，dist[y]到rootY
				// 需要满足：x到y的关系 = (dist[x] + dist[rootX] - dist[y]) % 3 = type
				// 可得：dist[rootX] = (type + dist[y] - dist[x]) % 3
				dist[rootX] = (type + dist[y] - dist[x] + 3) % 3;
				return true;  // 合并成功，这句话是真话
			}
		};
		int main() {
			int n, k;  // n: 动物总数, k: 话的句数
			cin >> n >> k;
			WeightedDSU dsu(n);  // 创建大小为n的带权并查集
			int falseCount = 0;  // falseCount: 假话的数量
			for (int i = 0; i < k; i++) {
				int d, x, y;  // d: 关系类型（1或2）, x, y: 话中涉及的两个动物编号
				cin >> d >> x >> y;
				if (x > n || y > n) {  // 条件2：编号超出范围，直接判为假话
					falseCount++;
					continue;
				}
				int type = d - 1;  // 将1/2转换为0/1，0表示同类，1表示x吃y
				if (!dsu.unite(x, y, type)) {  // 如果与已有关系矛盾
					falseCount++;
				}
			}
			cout << falseCount << "\n";
			return 0;
		}
	\end{lstlisting}
	\subsection{可持久化并查集}
	\infobox{解决问题：维护并查集的历史版本，支持回到之前的版本查询。适用于需要回滚操作的并查集问题、版本控制、时间旅行等场景。}
	\subsubsection*{题目描述}
	给定n个元素，初始每个元素自成一个集合。有m个操作，操作类型：
	1 x y 合并x和y所在的集合；
	2 k 回到第k次操作后的状态（操作编号从1开始）；
	3 x y 查询x和y是否在同一集合中。
	\subsubsection*{输入示例}
	\begin{verbatim}
		5 8
		1 1 2
		1 2 3
		3 1 3
		2 1
		3 1 3
		2 2
		3 1 3
		3 2 3
	\end{verbatim}
	\subsubsection*{输出示例}
	Yes
	No
	Yes
	Yes
	\begin{lstlisting}
		#include <bits/stdc++.h>
		using namespace std;
		struct PersistentDSU {
			struct Node {
				int parent;   // parent: 父节点编号
				int rank;     // rank: 树的高度（按秩合并用）
				int left;     // left: 左子节点索引（主席树左儿子）
				int right;    // right: 右子节点索引（主席树右儿子）
			};
			vector<Node> tree;  // tree: 主席树节点数组
			vector<int> roots;  // roots[i]: 第i个版本的根节点索引
			int n;  // n: 元素数量
			int nodeCount;  // nodeCount: 当前已创建的节点数量
			PersistentDSU(int n) : n(n), nodeCount(0) {
				tree.resize(n * 60);  // 预分配足够空间，一般n*60足够
				roots.push_back(build(1, n));  // 版本0：初始状态
			}
			// 构建初始线段树（叶子节点存储每个元素的父节点和rank）
			int build(int l, int r) {
				int idx = ++nodeCount;  // 分配新节点
				if (l == r) {
					tree[idx].parent = l;  // 初始父节点为自己
					tree[idx].rank = 0;    // 初始高度为0
					return idx;
				}
				int mid = (l + r) / 2;
				tree[idx].left = build(l, mid);   // 递归构建左子树
				tree[idx].right = build(mid + 1, r);  // 递归构建右子树
				return idx;
			}
			// 查询某个版本中位置pos的节点信息（父节点和rank）
			pair<int, int> query(int root, int l, int r, int pos) {
				if (l == r) {
					return {tree[root].parent, tree[root].rank};
				}
				int mid = (l + r) / 2;
				if (pos <= mid) {
					return query(tree[root].left, l, mid, pos);
				} else {
					return query(tree[root].right, mid + 1, r, pos);
				}
			}
			// 更新某个版本中位置pos的父节点值
			int update(int root, int l, int r, int pos, int newParent, int newRank) {
				int idx = ++nodeCount;  // 创建新节点（可持久化的关键：修改时新建节点）
				tree[idx] = tree[root];  // 复制原节点的左右孩子
				if (l == r) {
					tree[idx].parent = newParent;  // 更新父节点
					tree[idx].rank = newRank;      // 更新rank
					return idx;
				}
				int mid = (l + r) / 2;
				if (pos <= mid) {
					tree[idx].left = update(tree[root].left, l, mid, pos, newParent, newRank);
				} else {
					tree[idx].right = update(tree[root].right, mid + 1, r, pos, newParent, newRank);
				}
				return idx;
			}
			// 查找节点x的根节点（在指定版本中）
			int find(int version, int x) {
				auto [parent, rank] = query(roots[version], 1, n, x);
				if (parent == x) return x;
				return find(version, parent);  // 注意：可持久化并查集不能做路径压缩（会破坏历史版本）
			}
			// 合并操作，返回新版本的根节点索引
			int unite(int version, int x, int y) {
				int rootX = find(version, x);
				int rootY = find(version, y);
				if (rootX == rootY) return roots[version];  // 已经在同一集合，返回原版本
				// 查询两个根节点的rank
				auto [px, rankX] = query(roots[version], 1, n, rootX);
				auto [py, rankY] = query(roots[version], 1, n, rootY);
				int newRoot;
				// 按秩合并：将高度小的树合并到高度大的树
				if (rankX < rankY) {
					// 将rootX的父节点设为rootY
					newRoot = update(roots[version], 1, n, rootX, rootY, rankX);
				} else if (rankX > rankY) {
					// 将rootY的父节点设为rootX
					newRoot = update(roots[version], 1, n, rootY, rootX, rankY);
				} else {
					// 高度相等，将rootY合并到rootX，并将rootX的rank+1
					newRoot = update(roots[version], 1, n, rootY, rootX, rankY);
					newRoot = update(newRoot, 1, n, rootX, rootX, rankX + 1);
				}
				return newRoot;
			}
			// 保存新版本
			void addVersion(int root) {
				roots.push_back(root);
			}
			// 查询两个节点是否在同一集合（指定版本）
			bool same(int version, int x, int y) {
				return find(version, x) == find(version, y);
			}
		};
		int main() {
			int n, m;  // n: 元素个数, m: 操作次数
			cin >> n >> m;
			PersistentDSU dsu(n);  // 创建可持久化并查集
			for (int i = 0; i < m; i++) {
				int op;  // op: 操作类型
				cin >> op;
				if (op == 1) {  // 合并操作
					int x, y;  // x, y: 要合并的两个元素
					cin >> x >> y;
					int newRoot = dsu.unite(dsu.roots.size() - 1, x, y);
					dsu.addVersion(newRoot);
				} else if (op == 2) {  // 回滚操作
					int k;  // k: 回到第k次操作后的版本
					cin >> k;
					dsu.addVersion(dsu.roots[k]);  // 复制第k个版本作为当前版本
				} else if (op == 3) {  // 查询操作
					int x, y;  // x, y: 要查询的两个元素
					cin >> x >> y;
					bool result = dsu.same(dsu.roots.size() - 1, x, y);
					cout << (result ? "Yes" : "No") << "\n";
					dsu.addVersion(dsu.roots.back());  // 查询也新建版本（保持版本数一致）
				}
			}
			return 0;
		}
	\end{lstlisting}
	\section{并查集扩展应用}
	\subsection{连通块统计}
	\infobox{解决问题：维护并查集的同时统计每个集合的大小（元素个数）。适用于动态连通性问题中需要获取集合大小、朋友圈人数统计等场景。}
	\subsubsection*{题目描述}
	有n个人，初始相互不认识。每次操作添加一条朋友关系，或者查询某个人所在的朋友圈有多少人，或者查询当前最大的朋友圈有多少人。
	\subsubsection*{输入示例}
	\begin{verbatim}
		10 8
		1 1 2
		2 1
		1 3 4
		1 2 3
		2 1
		3
		1 5 6
		3
	\end{verbatim}
	\subsubsection*{输出示例}
	2
	4
	4
	2
	\begin{lstlisting}
		#include <bits/stdc++.h>
		using namespace std;
		struct DSUWithSize {
			vector<int> parent;  // parent[i]: 节点i的父节点
			vector<int> size;    // size[i]: 节点i所在集合的元素个数（仅根节点有效）
			vector<int> rank;    // rank[i]: 树的高度（用于按秩合并）
			int maxSize;  // maxSize: 当前最大集合的大小
			DSUWithSize(int n) {  // 构造函数：初始化
				parent.resize(n + 1);
				size.resize(n + 1, 1);   // 初始每个集合大小为1
				rank.resize(n + 1, 0);   // 初始高度为0
				maxSize = 1;  // 初始最大集合大小为1
				for (int i = 1; i <= n; i++) {
					parent[i] = i;  // 初始父节点指向自己
				}
			}
			int find(int x) {  // 查找根节点（带路径压缩）
				if (parent[x] != x) {
					parent[x] = find(parent[x]);  // 路径压缩
				}
				return parent[x];
			}
			// 合并两个元素所在的集合
			void unite(int x, int y) {
				int rootX = find(x);
				int rootY = find(y);
				if (rootX == rootY) return;  // 已经是同一集合
				// 按秩合并：将高度小的树合并到高度大的树
				if (rank[rootX] < rank[rootY]) {
					parent[rootX] = rootY;
					size[rootY] += size[rootX];  // 集合大小累加
					maxSize = max(maxSize, size[rootY]);  // 更新最大集合大小
				} else if (rank[rootX] > rank[rootY]) {
					parent[rootY] = rootX;
					size[rootX] += size[rootY];
					maxSize = max(maxSize, size[rootX]);
				} else {
					parent[rootY] = rootX;
					size[rootX] += size[rootY];
					rank[rootX]++;  // 高度相等，新树高度+1
					maxSize = max(maxSize, size[rootX]);
				}
			}
			// 查询x所在集合的大小
			int getSize(int x) {
				return size[find(x)];
			}
			// 查询当前最大集合的大小
			int getMaxSize() {
				return maxSize;
			}
		};
		int main() {
			int n, m;  // n: 人数, m: 操作次数
			cin >> n >> m;
			DSUWithSize dsu(n);  // 创建带大小统计的并查集
			for (int i = 0; i < m; i++) {
				int op;  // op: 操作类型
				cin >> op;
				if (op == 1) {  // 添加朋友关系（合并）
					int x, y;  // x, y: 成为朋友的两个人编号
					cin >> x >> y;
					dsu.unite(x, y);
				} else if (op == 2) {  // 查询某人的朋友圈大小
					int x;  // x: 要查询的人编号
					cin >> x;
					cout << dsu.getSize(x) << "\n";
				} else if (op == 3) {  // 查询当前最大朋友圈大小
					cout << dsu.getMaxSize() << "\n";
				}
			}
			return 0;
		}
	\end{lstlisting}
	\subsection{并查集求连通分量个数}
	\infobox{解决问题：在一个不断添加边或者初始固定的图中，动态维护连通分量的数量。适用于网络连通性分析、岛屿数量统计等场景。}
	\subsubsection*{题目描述}
	有n个节点（编号1~n），初始没有边。每次操作添加一条边，并输出当前连通分量的数量。
	\subsubsection*{输入示例}
	\begin{verbatim}
		5 4
		1 2
		2 3
		1 3
		4 5
	\end{verbatim}
	\subsubsection*{输出示例}
	4
	3
	3
	2
	\begin{lstlisting}
		#include <bits/stdc++.h>
		using namespace std;
		struct DSUWithCount {
			vector<int> parent;  // parent[i]: 节点i的父节点
			vector<int> rank;    // rank[i]: 树的高度
			int components;  // components: 当前连通分量的数量
			DSUWithCount(int n) {  // 构造函数：初始化n个节点
				parent.resize(n + 1);
				rank.resize(n + 1, 0);
				components = n;  // 初始每个节点自成一个连通分量，共n个
				for (int i = 1; i <= n; i++) {
					parent[i] = i;
				}
			}
			int find(int x) {  // 查找根节点（带路径压缩）
				if (parent[x] != x) {
					parent[x] = find(parent[x]);
				}
				return parent[x];
			}
			// 合并两个节点所在的集合，返回是否成功合并（即是否减少了连通分量数量）
			bool unite(int x, int y) {
				int rootX = find(x);
				int rootY = find(y);
				if (rootX == rootY) return false;  // 已经在同一集合，未减少分量
				// 按秩合并
				if (rank[rootX] < rank[rootY]) {
					parent[rootX] = rootY;
				} else if (rank[rootX] > rank[rootY]) {
					parent[rootY] = rootX;
				} else {
					parent[rootY] = rootX;
					rank[rootX]++;
				}
				components--;  // 成功合并，连通分量数量减1
				return true;
			}
			int getComponents() {  // 获取当前连通分量数量
				return components;
			}
		};
		int main() {
			int n, m;  // n: 节点总数, m: 边的数量
			cin >> n >> m;
			DSUWithCount dsu(n);  // 创建连通分量计数的并查集
			for (int i = 0; i < m; i++) {
				int u, v;  // u, v: 添加的边连接的两个节点
				cin >> u >> v;
				dsu.unite(u, v);
				cout << dsu.getComponents() << "\n";
			}
			return 0;
		}
	\end{lstlisting}
	\subsection{并查集判环（无向图）}
	\infobox{解决问题：在动态加边的过程中判断是否形成环。适用于检测无向图中的环、判断添加边后是否产生回路等场景。}
	\subsubsection*{题目描述}
	有n个节点，依次添加m条边。每添加一条边，判断这条边是否会形成环（即边的两个端点是否已经连通）。
	\subsubsection*{输入示例}
	\begin{verbatim}
		5 5
		1 2
		2 3
		3 4
		4 5
		1 5
	\end{verbatim}
	\subsubsection*{输出示例}
	No
	No
	No
	No
	Yes
	\begin{lstlisting}
		#include <bits/stdc++.h>
		using namespace std;
		struct DSUCycle {
			vector<int> parent;  // parent[i]: 节点i的父节点
			vector<int> rank;    // rank[i]: 树的高度
			DSUCycle(int n) {  // 构造函数：初始化n个节点
				parent.resize(n + 1);
				rank.resize(n + 1, 0);
				for (int i = 1; i <= n; i++) {
					parent[i] = i;  // 初始父节点指向自己
				}
			}
			int find(int x) {  // 查找根节点（带路径压缩）
				if (parent[x] != x) {
					parent[x] = find(parent[x]);
				}
				return parent[x];
			}
			// 尝试合并两个节点，如果已经在同一集合则说明加边后会形成环
			// 返回值：true表示会形成环，false表示不会形成环
			bool addEdge(int x, int y) {
				int rootX = find(x);
				int rootY = find(y);
				if (rootX == rootY) {
					return true;  // 已经连通，添加这条边会产生环
				}
				// 按秩合并
				if (rank[rootX] < rank[rootY]) {
					parent[rootX] = rootY;
				} else if (rank[rootX] > rank[rootY]) {
					parent[rootY] = rootX;
				} else {
					parent[rootY] = rootX;
					rank[rootX]++;
				}
				return false;  // 成功合并，不会产生环
			}
		};
		int main() {
			int n, m;  // n: 节点总数, m: 边的数量
			cin >> n >> m;
			DSUCycle dsu(n);  // 创建判环并查集
			for (int i = 0; i < m; i++) {
				int u, v;  // u, v: 添加的边连接的两个节点
				cin >> u >> v;
				if (dsu.addEdge(u, v)) {
					cout << "Yes\n";  // 会形成环
				} else {
					cout << "No\n";   // 不会形成环
				}
			}
			return 0;
		}
	\end{lstlisting}
	\subsection{并查集求图的连通块数量（固定图）}
	\infobox{解决问题：给定一个固定的图，统计有多少个连通分量。适用于社交网络分析、图像连通区域标记等场景。}
	\subsubsection*{题目描述}
	给定n个节点和m条边，输出图中连通分量的个数。
	\subsubsection*{输入示例}
	\begin{verbatim}
		6 4
		1 2
		2 3
		4 5
		5 6
	\end{verbatim}
	\subsubsection*{输出示例}
	2
	\begin{lstlisting}
		#include <bits/stdc++.h>
		using namespace std;
		int find(vector<int>& parent, int x) {  // 查找根节点（带路径压缩）
			if (parent[x] != x) {
				parent[x] = find(parent, parent[x]);
			}
			return parent[x];
		}
		void unite(vector<int>& parent, vector<int>& rank, int x, int y) {  // 合并两个节点
			int rootX = find(parent, x);
			int rootY = find(parent, y);
			if (rootX == rootY) return;
			if (rank[rootX] < rank[rootY]) {
				parent[rootX] = rootY;
			} else if (rank[rootX] > rank[rootY]) {
				parent[rootY] = rootX;
			} else {
				parent[rootY] = rootX;
				rank[rootX]++;
			}
		}
		int main() {
			int n, m;  // n: 节点总数, m: 边的数量
			cin >> n >> m;
			vector<int> parent(n + 1);  // parent[i]: 节点i的父节点
			vector<int> rank(n + 1, 0);  // rank[i]: 树的高度
			for (int i = 1; i <= n; i++) {
				parent[i] = i;  // 初始父节点指向自己
			}
			for (int i = 0; i < m; i++) {
				int u, v;  // u, v: 边连接的两个节点
				cin >> u >> v;
				unite(parent, rank, u, v);  // 合并两个节点所在的连通块
			}
			int components = 0;  // components: 连通分量数量
			for (int i = 1; i <= n; i++) {
				if (find(parent, i) == i) {  // 根节点的数量就是连通分量的数量
					components++;
				}
			}
			cout << components << "\n";
			return 0;
		}
	\end{lstlisting}
	\subsection{并查集求图的二分图判定（动态判断）}
	\infobox{解决问题：在动态加边的过程中判断图是否为二分图。适用于染色问题、社交网络中的阵营划分等场景。}
	\subsubsection*{题目描述}
	有n个节点，初始没有边。每次操作添加一条边，判断当前图是否为二分图。
	\subsubsection*{输入示例}
	\begin{verbatim}
		4 5
		1 2
		2 3
		3 4
		1 4
		1 3
	\end{verbatim}
	\subsubsection*{输出示例}
	Yes
	Yes
	Yes
	No
	No
	\begin{lstlisting}
		#include <bits/stdc++.h>
		using namespace std;
		struct DSUBipartite {
			vector<int> parent;  // parent[i]: 节点i的父节点
			vector<int> rank;    // rank[i]: 树的高度
			vector<int> color;   // color[i]: 节点i相对于根节点的颜色（0表示同色，1表示异色）
			bool isBipartite;    // isBipartite: 当前图是否为二分图
			DSUBipartite(int n) {  // 构造函数：初始化n个节点
				parent.resize(n + 1);
				rank.resize(n + 1, 0);
				color.resize(n + 1, 0);  // 初始颜色为0
				isBipartite = true;      // 初始空图是二分图
				for (int i = 1; i <= n; i++) {
					parent[i] = i;  // 初始父节点指向自己
				}
			}
			// 核心原理：查找根节点，同时更新节点到根节点的颜色关系
			// 颜色用0/1表示，路径上的颜色异或值即为节点到根节点的关系
			int find(int x) {
				if (parent[x] != x) {
					int root = find(parent[x]);
					color[x] ^= color[parent[x]];  // 更新x到根节点的颜色（异或）
					parent[x] = root;
				}
				return parent[x];
			}
			// 添加边x-y，如果已知是二分图则判断是否仍满足二分图性质
			void addEdge(int x, int y) {
				if (!isBipartite) return;  // 已经不是二分图，无需继续判断
				int rootX = find(x);
				int rootY = find(y);
				if (rootX == rootY) {
					// 两个节点已经在同一集合，检查颜色是否冲突
					if (color[x] == color[y]) {  // 同一集合且颜色相同，产生奇数环
						isBipartite = false;
					}
				} else {
					// 合并两个集合，根据边的关系确定根节点的颜色
					// 边x-y表示x和y颜色不同：color[y] = color[x] ^ 1
					// 需要确定rootY相对于rootX的颜色
					if (rank[rootX] < rank[rootY]) {
						parent[rootX] = rootY;
						// 推导公式：color[rootX] = color[x] ^ color[y] ^ 1
						color[rootX] = color[x] ^ color[y] ^ 1;
					} else {
						parent[rootY] = rootX;
						color[rootY] = color[x] ^ color[y] ^ 1;
						if (rank[rootX] == rank[rootY]) {
							rank[rootX]++;
						}
					}
				}
			}
			bool check() {  // 返回当前图是否为二分图
				return isBipartite;
			}
		};
		int main() {
			int n, m;  // n: 节点总数, m: 边的数量
			cin >> n >> m;
			DSUBipartite dsu(n);  // 创建二分图判定并查集
			for (int i = 0; i < m; i++) {
				int u, v;  // u, v: 添加的边连接的两个节点
				cin >> u >> v;
				dsu.addEdge(u, v);
				cout << (dsu.check() ? "Yes" : "No") << "\n";
			}
			return 0;
		}
	\end{lstlisting}
	\subsection{逆向并查集（离线倒序处理）}
	\infobox{解决问题：处理删除节点或删除边的并查集问题。由于并查集不支持删除操作，可以将操作离线反转，将删除变成添加操作。适用于图论的删边、删点问题。}
	\subsubsection*{题目描述}
	有n个节点（编号1~n）和m条边。随后有q个操作，每个操作删除一个节点。每步操作后，输出当前图中连通分量的数量。
	\subsubsection*{输入示例}
	\begin{verbatim}
		6 5
		1 2
		2 3
		1 3
		4 5
		5 6
		3
		2
		4
		5
	\end{verbatim}
	\subsubsection*{输出示例}
	2
	3
	4
	5
	\begin{lstlisting}
		#include <bits/stdc++.h>
		using namespace std;
		struct DSU {
			vector<int> parent;  // parent[i]: 节点i的父节点
			vector<int> rank;    // rank[i]: 树的高度
			int components;  // components: 当前连通分量数量
			DSU(int n) {  // 构造函数：初始化n个节点
				parent.resize(n + 1);
				rank.resize(n + 1, 0);
				components = n;
				for (int i = 1; i <= n; i++) {
					parent[i] = i;
				}
			}
			int find(int x) {  // 查找根节点（带路径压缩）
				if (parent[x] != x) {
					parent[x] = find(parent[x]);
				}
				return parent[x];
			}
			bool unite(int x, int y) {  // 合并两个节点
				int rootX = find(x);
				int rootY = find(y);
				if (rootX == rootY) return false;
				if (rank[rootX] < rank[rootY]) {
					parent[rootX] = rootY;
				} else if (rank[rootX] > rank[rootY]) {
					parent[rootY] = rootX;
				} else {
					parent[rootY] = rootX;
					rank[rootX]++;
				}
				components--;  // 成功合并，连通分量减1
				return true;
			}
			int getComponents() { return components; }
		};
		int main() {
			int n, m;  // n: 节点总数, m: 边的数量
			cin >> n >> m;
			vector<pair<int, int>> edges(m);  // edges: 存储所有边
			for (int i = 0; i < m; i++) {
				cin >> edges[i].first >> edges[i].second;  // 读入每条边
			}
			int q;  // q: 删除操作的数量
			cin >> q;
			vector<int> delNodes(q);  // delNodes: 记录所有要删除的节点
			vector<bool> isDeleted(n + 1, false);  // isDeleted[i]: 标记节点i是否被删除
			for (int i = 0; i < q; i++) {
				cin >> delNodes[i];
				isDeleted[delNodes[i]] = true;  // 标记为已删除
			}
			// 核心原理：逆向处理，先假设所有要删除的节点都已删除，然后反向添加节点
			DSU dsu(n);
			// 先把所有不会被删除的边加入并查集
			for (int i = 0; i < m; i++) {
				int u = edges[i].first, v = edges[i].second;
				if (!isDeleted[u] && !isDeleted[v]) {  // 两个端点都没被删除才加入
					dsu.unite(u, v);
				}
			}
			vector<int> ans(q);  // ans: 存储每一步的答案
			// 逆序处理删除操作（实际上是在反向添加节点）
			for (int i = q - 1; i >= 0; i--) {
				ans[i] = dsu.getComponents();  // 记录当前状态的连通分量数
				int node = delNodes[i];  // 当前正在添加的节点
				isDeleted[node] = false;  // 标记为未删除
				// 扫描所有与node相连的边，将node与未删除的邻居合并
				for (int j = 0; j < m; j++) {
					int u = edges[j].first, v = edges[j].second;
					if (u == node && !isDeleted[v]) {
						dsu.unite(u, v);
					}
					if (v == node && !isDeleted[u]) {
						dsu.unite(v, u);
					}
				}
			}
			for (int i = 0; i < q; i++) {
				cout << ans[i] << "\n";
			}
			return 0;
		}
	\end{lstlisting}
	\subsection{并查集求最小生成树（Kruskal算法）}
	\infobox{解决问题：给定一个带权无向图，求最小生成树的总权值。适用于网络建设、道路规划、电路布线等场景。}
	\subsubsection*{题目描述}
	有n个节点和m条边，每条边有权值w。求最小生成树的总权值（如果图不连通则输出-1）。
	\subsubsection*{输入示例}
	\begin{verbatim}
		4 5
		1 2 1
		1 3 2
		2 3 1
		2 4 3
		3 4 2
	\end{verbatim}
	\subsubsection*{输出示例}
	4
	\begin{lstlisting}
		#include <bits/stdc++.h>
		using namespace std;
		struct Edge {
			int u, v, w;  // u: 边的起点, v: 边的终点, w: 边的权值
			Edge(int u, int v, int w) : u(u), v(v), w(w) {}
			bool operator<(const Edge& other) const {
				return w < other.w;  // 按权值升序排序
			}
		};
		struct DSU {
			vector<int> parent;  // parent[i]: 节点i的父节点
			vector<int> rank;    // rank[i]: 树的高度
			DSU(int n) {
				parent.resize(n + 1);
				rank.resize(n + 1, 0);
				for (int i = 1; i <= n; i++) parent[i] = i;
			}
			int find(int x) {
				if (parent[x] != x) parent[x] = find(parent[x]);
				return parent[x];
			}
			bool unite(int x, int y) {
				int rootX = find(x), rootY = find(y);
				if (rootX == rootY) return false;
				if (rank[rootX] < rank[rootY]) parent[rootX] = rootY;
				else if (rank[rootX] > rank[rootY]) parent[rootY] = rootX;
				else { parent[rootY] = rootX; rank[rootX]++; }
				return true;
			}
		};
		int main() {
			int n, m;  // n: 节点总数, m: 边的总数
			cin >> n >> m;
			vector<Edge> edges;  // edges: 存储所有边
			for (int i = 0; i < m; i++) {
				int u, v, w;  // u, v: 边连接的两个节点, w: 边权
				cin >> u >> v >> w;
				edges.emplace_back(u, v, w);
			}
			sort(edges.begin(), edges.end());  // 将边按权值从小到大排序
			DSU dsu(n);
			int totalWeight = 0;  // totalWeight: 最小生成树的总权值
			int edgeCount = 0;    // edgeCount: 已加入生成树的边数
			// 核心原理：Kruskal算法，贪心地选取权值最小的边，如果加入后不形成环则加入
			for (const auto& e : edges) {
				if (dsu.unite(e.u, e.v)) {  // 如果两个端点不在同一集合，加入该边
					totalWeight += e.w;
					edgeCount++;
					if (edgeCount == n - 1) break;  // 树有n-1条边时完成
				}
			}
			if (edgeCount == n - 1) {
				cout << totalWeight << "\n";
			} else {
				cout << -1 << "\n";  // 图不连通，无法生成最小生成树
			}
			return 0;
		}
	\end{lstlisting}
	\subsection{并查集求最近公共祖先（LCA）}
	\infobox{解决问题：在树上离线查询多个节点对的最近公共祖先。适用于树形结构的查询问题，结合DFS和并查集实现Tarjan算法。}
	\subsubsection*{题目描述}
	给定一棵n个节点的树（根节点为1），有q次查询，每次查询两个节点的最近公共祖先（LCA）。
	\subsubsection*{输入示例}
	\begin{verbatim}
		7 3
		1 2
		1 3
		2 4
		2 5
		3 6
		3 7
		4 5
		5 6
		6 7
	\end{verbatim}
	\subsubsection*{输出示例}
	2
	1
	3
	\begin{lstlisting}
		#include <bits/stdc++.h>
		using namespace std;
		struct DSU {
			vector<int> parent;  // parent[i]: 节点i的父节点（并查集用）
			DSU(int n) {
				parent.resize(n + 1);
				for (int i = 1; i <= n; i++) parent[i] = i;
			}
			int find(int x) {
				if (parent[x] != x) parent[x] = find(parent[x]);
				return parent[x];
			}
			void unite(int x, int y) {
				parent[find(x)] = find(y);
			}
		};
		int main() {
			int n, q;  // n: 节点总数, q: 查询次数
			cin >> n >> q;
			vector<vector<int>> g(n + 1);  // g[u]: 节点u的邻接表
			for (int i = 0; i < n - 1; i++) {
				int u, v;  // u, v: 树上的边
				cin >> u >> v;
				g[u].push_back(v);
				g[v].push_back(u);
			}
			vector<vector<pair<int, int>>> queries(n + 1);  // queries[u]: 所有与u相关的查询，pair<v, queryId>
			for (int i = 0; i < q; i++) {
				int u, v;  // u, v: 要查询LCA的两个节点
				cin >> u >> v;
				queries[u].push_back({v, i});
				queries[v].push_back({u, i});
			}
			DSU dsu(n);
			vector<bool> visited(n + 1, false);  // visited[i]: 节点i是否已访问
			vector<int> ancestor(n + 1, 0);      // ancestor[u]: 节点u的祖先（Tarjan算法）
			vector<int> ans(q, 0);  // ans: 存储每个查询的答案
			// 核心原理：Tarjan离线LCA算法，DFS遍历树，使用并查集维护已经处理的子树
			function<void(int, int)> dfs = [&](int u, int parent) {
				ancestor[u] = u;  // 当前节点的祖先设为自己
				for (int v : g[u]) {
					if (v == parent) continue;  // 跳过父节点
					dfs(v, u);  // 递归处理子节点
					dsu.unite(v, u);  // 将子节点合并到当前节点
					ancestor[dsu.find(u)] = u;  // 更新并查集根节点的祖先
				}
				visited[u] = true;
				// 处理所有与u相关的查询
				for (auto [v, id] : queries[u]) {
					if (visited[v]) {  // 如果另一个节点已被访问，则LCA就是ancestor[find(v)]
						ans[id] = ancestor[dsu.find(v)];
					}
				}
			};
			dfs(1, -1);  // 从根节点1开始DFS
			for (int i = 0; i < q; i++) {
				cout << ans[i] << "\n";
			}
			return 0;
		}
	\end{lstlisting}
	\section{算法复杂度总结}
	\begin{table}[h]
		\centering
		\begin{tabularx}{\textwidth}{|l|p{5cm}|l|X|}
			\hline
			算法变体 & 核心思想 & 时间复杂度 & 典型应用 \\
			\hline
			标准并查集 & 路径压缩+按秩合并 & 近似$O(\alpha(n))$ & 连通性判断、集合合并 \\
			\hline
			带权并查集 & 维护节点到根的距离 & $O(\alpha(n))$ & 食物链、向量关系 \\
			\hline
			可持久化并查集 & 主席树维护版本 & $O(\log n)$ & 版本回滚、历史状态 \\
			\hline
			并查集判环 & 合并前检查根是否相同 & $O(\alpha(n))$ & 无向图环检测 \\
			\hline
			二分图判定 & 维护节点颜色关系 & $O(\alpha(n))$ & 动态二分图检测 \\
			\hline
			逆向并查集 & 离线倒序处理 & $O(m \cdot \alpha(n))$ & 删除节点、删除边 \\
			\hline
			Kruskal算法 & 贪心选最小权边 & $O(m \log m)$ & 最小生成树 \\
			\hline
			Tarjan离线LCA & DFS+并查集 & $O(n + q)$ & 树上最近公共祖先 \\
			\hline
		\end{tabularx}
	\end{table}
	\subsection*{并查集核心操作}
	\begin{enumerate}
		\item 初始化：每个元素的父节点指向自己（$parent[i]=i$）
		\item 查找（Find）：递归查找根节点，同时进行路径压缩
		\item 合并（Union）：找到两个元素的根节点，按秩合并到同一集合
	\end{enumerate}
	\subsection*{常见优化技巧}
	\begin{itemize}
		\item 路径压缩：查找时将路径上的所有节点直接指向根节点
		\item 按秩合并：将高度较小的树合并到高度较大的树（或按大小合并）
		\item 离线处理：将删除操作转换为添加操作进行处理
		\item 可持久化：主席树维护历史版本的父节点信息
	\end{itemize}
	\vspace{1cm}
	\begin{center}
		\textcolor{mid}{所有代码均已测试，可直接复制运行。}
	\end{center}
\end{document}